% Generated by tools/build_papers.py from the authoritative Markdown master.
\documentclass[11pt,a4paper]{article}
\usepackage{fontspec}
\usepackage{amsmath,amssymb,mathtools}
\usepackage[margin=25mm]{geometry}
\usepackage{graphicx,booktabs,longtable,array,calc,enumitem}
\usepackage{microtype}
\usepackage{needspace,xurl}
\usepackage[unicode,hidelinks]{hyperref}
\usepackage{bookmark}
\defaultfontfeatures{Ligatures=TeX}
\setmainfont{Latin Modern Roman}
\setsansfont{Latin Modern Sans}
\setmonofont{Latin Modern Mono}[Scale=0.91]
\setlength{\emergencystretch}{5em}
\setlength{\parindent}{1.2em}
\setlength{\parskip}{2pt plus 1pt}
\linespread{1.05}
\raggedbottom
\clubpenalty=10000
\widowpenalty=10000
\displaywidowpenalty=10000
\allowdisplaybreaks[2]
\setcounter{secnumdepth}{3}
\setcounter{tocdepth}{2}
\providecommand{\tightlist}{\setlength{\itemsep}{0pt}\setlength{\parskip}{0pt}}
\setlist{itemsep=2pt,parsep=0pt,topsep=5pt}
\newcommand{\passthrough}[1]{#1}
\makeatletter
\newcommand{\paperstatement}[3]{\par\addvspace{6pt}\Needspace{4\baselineskip}\phantomsection\def\@currentlabel{#2}\label{#3}\noindent\textbf{#1}\enspace}
\makeatother
\newcommand{\paperref}[2]{\hyperref[#1]{#2}}
\makeatletter
\renewcommand{\maketitle}{\begin{center}{\LARGE\bfseries\@title\par}\vspace{1em}{\normalsize\@author\par}\vspace{0.5em}{\small\@date\par}\end{center}\vspace{1em}}
\makeatother
\title{Constant-factor bounds for Hamiltonian paths in tournaments}
\author{Xingchen Liu and Xiangyu Ye\\Independent Researchers\\lxc-em5158@outlook.com}
\date{8 October 2026; revised 11 October 2026}
\hypersetup{pdftitle={Constant-factor bounds for Hamiltonian paths in tournaments},pdfauthor={Xingchen Liu and Xiangyu Ye},pdfsubject={Hamiltonian paths in tournaments},bookmarksnumbered=true}
\begin{document}
\maketitle
\section*{Abstract}
\phantomsection\addcontentsline{toc}{section}{Abstract}\label{sec-abstract}

Let \(P(n)\) be the maximum number of directed Hamiltonian paths in an \(n\)-vertex tournament, and let \(\mu_n=n!/2^{n-1}\) be the random-tournament expectation. We prove
\[(L-O(n^{-1}))\mu_n\le P(n)\le(C_*+O(n^{-1}))\mu_n,\]

where
\[L=\frac{\cosh1}{\cos1}=2.855957892565\ldots,
\qquad
C_*=\frac{3\pi^4+4\pi^2-32}{\pi^4+4\pi^2-32}
=2.857401177672\ldots.\]

The upper bound holds for every tournament. Its proof starts from a positive determinant--permanent identity for the path count. A uniform permanent approximation and quantitative matrix scaling convert this identity into a spectral estimate, while degree variance and a multiplicative score penalty control irregular tournaments. A skew spectral-radius bound then gives \(C_*\) by convex optimization. Carousel tournaments of both parities give the lower bound. The relative gap between the two leading constants is about \(0.050536\%\). A companion Lean development proves the main theorem and the uniform small-score path approximation.

Keywords: tournament; directed Hamiltonian path; permanent; matrix scaling; skew spectrum.

\Needspace{9\baselineskip}
\section{Introduction and main result}\label{sec-1}

How many directed Hamiltonian paths can a tournament contain? For an \(n\)-vertex tournament \(T\), let \(H(T)\) count the orderings \((v_1,\ldots,v_n)\) for which \(v_j\to v_{j+1}\) for every \(j<n\). Paths are counted as directed vertex sequences. Define
\[P(n)=\max_{|V(T)|=n}H(T),\qquad \mu_n=\frac{n!}{2^{n-1}}.\]

Each ordering forms a directed path in a uniformly random tournament with probability \(2^{-(n-1)}\). Thus \(\mathbb E H(T)=\mu_n\) and \(P(n)\ge\mu_n\). The question is how much a suitable orientation can improve on this expectation.

Alon \paperref{bib-2}{[2]} proved \(P(n)=O(n^{3/2})\mu_n\) using permanents. Friedgut and Kahn\textquotesingle s Hamiltonian-cycle bound \paperref{bib-7}{[7]}, combined with the path-to-cycle construction in \paperref{bib-2}{[2, Proposition 2.5]}, gives \(P(n)=O(n^{3/2-\xi})\mu_n\), where \(\xi\approx0.2507\). Wormald \paperref{bib-6}{[6, Theorem 5]} obtained \(P(n)>2.85588\mu_n\) for infinitely many \(n\) and conjectured a limiting ratio approximately \(2.855958\). Our result places the maximum path count between two close constant multiples of \(\mu_n\).

\Needspace{12\baselineskip}\paperstatement{Theorem 1.1 (Main result).}{1.1}{thm-1.1} There are absolute constants \(K\ge0\) and \(n_0\ge2\) such that, for every integer \(n\ge n_0\),
\[(L-K/n)\mu_n\le P(n)\le(C_*+K/n)\mu_n,\]

where
\[L=\frac{\cosh1}{\cos1},\qquad
a_*=\frac4{\pi^2},\qquad
C_*=\frac{1+a_*}{1-a_*}\frac{3/2-a_*}{1/2+a_*}.\]

The main task is the upper bound for arbitrary tournaments. Brégman\textquotesingle s inequality controls a permanent from its row sums, but even for regular tournaments its bound leaves a factor of order \(\sqrt n\) at the natural permanent scale \(n!/2^n\). Removing this loss requires the relations between different rows imposed by the tournament orientation.

We use the determinant--permanent path identity of Irving and Omar \paperref{bib-1}{[1]}. Its summands are nonnegative, and only small deleted vertex sets contribute appreciably. This reduces the counting problem to accurate estimates for permanents of large submatrices. For nearly regular tournaments, these estimates yield a spectral factor \(\rho_n(S)\), where \(S=A-A^{\mathsf T}\) is the skew adjacency matrix. \paperref{sec-2.3}{Section 2.3} explains its origin before the technical proofs.

Three points require quantitative control. First, the permanent approximation must hold uniformly for a whole class of matrices, rather than only for each fixed coefficient degree. Second, the actual adjacency submatrices must be scaled to have equal row and column sums; different deleted row and column sets arise when exceptional vertices are separated. Third, arbitrary tournaments must be reduced to a class in which this scaling is available. In that class, the restored scaling factors contain a negative score exponent that absorbs the approximation errors.

We present the counting argument first. \paperref{sec-2}{Section 2} establishes the exact identities and the spectral bound. \paperref{sec-3}{Section 3} proves the upper bound using a precisely stated permanent estimate, and \paperref{sec-4}{Section 4} proves the small-score formula and the carousel lower bound. \paperref{sec-5}{Sections 5} and 6 then prove the permanent and scaling estimates. The permanent argument builds on the determinantal viewpoint of McCullagh \paperref{bib-5}{[5]} and Gaussian permanent tools used by Han and Niles-Weed \paperref{bib-3}{[3]}; the skew spectral-radius comparison is given in Deng, Li, Shader and So \paperref{bib-4}{[4]}. We include the versions and proofs needed here. \paperref{sec-A}{Appendix A} records the correspondence with the companion formalization.

\Needspace{9\baselineskip}
\section{Counting identities and the spectral factor}\label{sec-2}

\Needspace{7\baselineskip}
\subsection{Matrices, scores and spectral factors}\label{sec-2.1}

The adjacency matrix \(A\) has zero diagonal and \(A_{ij}=1\) precisely when \(i\to j\). Put
\[S=A-A^{\mathsf T},\qquad 2A=J-I+S,\qquad s=S\mathbf1.\]

Here \(I\) is the identity, \(J=\mathbf1\mathbf1^{\mathsf T}\) and \(\mathbf1\) is the all-ones column vector. If \(d_i^+\) is the outdegree, then \(s_i=2d_i^+-(n-1)\) and \(\sum_i s_i=0\). We use
\[\mathcal E=\|s\|_2^2,\qquad d=\|s\|_\infty,\qquad
a=\frac{s}{n-1},\qquad \tau=\|a\|_2^2.\]

The normalized scores \(a\) and \(\tau\) are defined for \(n\ge2\). A regular odd-order tournament has \(s=0\), and a balanced even-order tournament has \(s_i\in\{-1,1\}\). Thus \(d\) measures the largest vertex imbalance, while \(\tau\) measures its total squared size after normalization.

For a vertex set \(U\), the notation \(A[U]\) denotes a principal submatrix. For possibly different row and column sets \(R,C\), write \(A[R,C]\). The determinant and permanent of the empty matrix are both one. The permanent is
\[\operatorname{per}M=\sum_{\pi\in\mathfrak S_m}\prod_{i=1}^m M_{i,\pi(i)}.\]

For a real matrix \(M\), its singular values are the square roots of the eigenvalues of \(M^{\mathsf T}M\). Their maximum is \(\|M\|_{\rm op}\), and the sum of their squares is \(\|M\|_F^2=\sum_{i,j}|M_{ij}|^2\). For the skew matrix \(S\), write \(\pm i\lambda_j\), with \(\lambda_j>0\), for the nonzero eigenvalue pairs, and add zero frequencies when convenient. Each \(\lambda_j\) occurs twice among the singular values. Define
\[x_j=\frac{\lambda_j^2}{n^2},\qquad
D_n(S)=\det(I-S^{\mathsf T}S/n^2)^{-1/2},\]
\nopagebreak[3]\[\rho_n(S)=D_n(S)\det(I+S/n)=\prod_j\frac{1+x_j}{1-x_j}.\]

Skew singular values occur in pairs, and \(\|S\|_F^2=n(n-1)\), whence
\[\sum_jx_j=\frac{n-1}{2n},\qquad \|S/n\|_{\rm op}^2\le\frac{n-1}{2n}<\frac12.\]

Thus \(D_n\) and \(\rho_n\) are positive and bounded by absolute constants. In particular,
\[D_n(S)=\det(I+iS/n)^{-1}.\]

\Needspace{7\baselineskip}
\subsection{The positive path convolution}\label{sec-2.2}

\paperstatement{Lemma 2.1 (Path convolution).}{2.1}{thm-2.1} For every tournament,
\[H(T)=\sum_{U\subseteq V(T)}\det(I+A[U])\operatorname{per}A[U^c].\]

The permanent of an adjacency matrix counts directed cycle covers. This identity converts those counts on complementary vertex sets into a path count. It is the tournament specialization of Irving and Omar\textquotesingle s Proposition 2 \paperref{bib-1}{[1]}, with complement adjacency matrix \(\overline A=J-A=I+A^{\mathsf T}\); the complement includes its diagonal entries.

Here is also a generating-function verification. Put \(X=\operatorname{diag}(z_1,\ldots,z_n)\). The generating series of all walks, with each vertex occurrence carrying its variable, is
\[1+\mathbf1^{\mathsf T}(I-XA)^{-1}X\mathbf1
=\frac{\det(I+X\overline A)}{\det(I-XA)}.\]

The equality is the rank-one determinant lemma. Taking the coefficient of \(z_1\cdots z_n\) selects exactly the Hamiltonian paths. For the denominator, the formal identity \(\det(I-XA)^{-1}=\exp(\sum_{r\ge1}\operatorname{tr}((XA)^r)/r)\) shows that a squarefree coefficient counts disjoint directed cycle covers, hence a permanent. Combining this with the principal-minor expansion of the numerator gives
\[H(T)=\sum_U\det\overline A[U]\operatorname{per}A[U^c].\]

Using \(\overline A[U]=I+A[U]^{\mathsf T}\) proves the assertion.

If \(k=|U|\), the symmetric part of \(I+A[U]\) is \((I+J)/2\), which is positive definite. Every real eigenvalue is therefore positive, and nonreal eigenvalues occur in conjugate pairs, so its determinant is positive. The squared row norms are \(1+d_i^+(T[U])\), with mean \((k+1)/2\). Hadamard\textquotesingle s inequality followed by arithmetic--geometric mean therefore gives
\[0<\det(I+A[U])\le h_k,\qquad h_k=\left(\frac{k+1}{2}\right)^{k/2}.\]

Set \(h_0=1\) and
\[w_k=\frac{2^k h_k}{k!}.\]

For every fixed \(c>0\) and fixed nonnegative integer \(r\),
\[\sum_{k\ge0}k^r c^k w_k<\infty.\]

Indeed, Stirling\textquotesingle s formula gives \(\log(c^kw_k)=-(k/2)\log k+O_c(k)\). With
\[k_n=\left\lceil\frac{4\log n}{\log\log n}\right\rceil\]

for sufficiently large \(n\), it follows that \(\sum_{k>k_n}c^kw_k=n^{-2+o(1)}\).

Brégman\textquotesingle s bound in the form \paperref{bib-2}{[2, Lemma 2.1]}, together with degree balancing \paperref{bib-2}{[2, Corollary 2.3]} and Stirling\textquotesingle s formula, gives a universal constant \(C\) with
\[\operatorname{per}A_T\le C\sqrt{m+1}\,\frac{m!}{2^m}\]

for every tournament of order \(m\), including \(m=0\) after enlarging \(C\). \paperref{thm-3.2}{Lemma 3.2} proves a variance-penalized version of this bound. Consequently, the contribution to the path convolution from \(|U|>k_n\), divided by \(\mu_n\), is at most
\[\frac C2\sqrt{n+1}\sum_{k>k_n}w_k
=n^{-3/2+o(1)}=o(n^{-1}).\]

This bound is uniform in the tournament. We call \(|U|\le k_n\) the short terms. Their accurate estimation is the only permanent problem left by the convolution.

\Needspace{7\baselineskip}
\subsection{Why the spectral factor appears}\label{sec-2.3}

The following calculation explains the quantity to be bounded. Suppose for the moment that the maximum score satisfies \(d=o(\sqrt n)\). \paperref{thm-6.4}{Lemma 6.4} will show that, for a short deletion \(|U|=k\),
\[\operatorname{per}A[U^c]
=e^{-1}D_n(S)\frac{(n-k)!}{2^{n-k}}
\left(1+O\left(\frac{(d+k+1)^2}{n}\right)\right).\]

The normalization satisfies \(((n-k)!/2^{n-k})/\mu_n=2^{k-1}/(n)_k\). Thus inserting the leading term in \paperref{thm-2.1}{Lemma 2.1} gives the normalized sum
\[\frac{e^{-1}D_n(S)}2
\sum_{|U|\le k_n}\frac{2^{|U|}}{(n)_{|U|}}\det(I+A[U]),\]

where \((n)_k=n(n-1)\cdots(n-k+1)\). Replacing \((n)_k\) by \(n^k\) in the short sum and adding the negligible tail produces errors estimated in \paperref{sec-4.1}{Section 4.1}. The resulting full sum has the exact principal-minor identity
\[\sum_U(2/n)^{|U|}\det(I+A[U])
=\det\left(I+\frac2n(I+A)\right).\]

Since \(2A=J-I+S\), the rank-one term \(J\) contributes asymptotically a factor two, and the scalar term contributes \(e\). More precisely, the small-score assumption gives
\[\det\left(I+\frac2n(I+A)\right)
=2e\det(I+S/n)\left(1+O(n^{-1}+d^2/n^2)\right).\]

The \(e^{-1}\) in the permanent estimate comes from restoring \((1-1/n)^{n-k}\), which accounts for the zero diagonal. The factor \(1/2\) comes from the path normalization above. Both cancel against the generating determinant\textquotesingle s \(2e\), leaving \(D_n(S)\det(I+S/n)=\rho_n(S)\). The two determinants have separate origins: the first approximates cycle covers, and the second sums the path-convolution weights. \paperref{thm-4.1}{Theorem 4.1} supplies the error summation and proves
\[\frac{H(T)}{\mu_n}=\rho_n(S)+O((d+1)^2/n)
\qquad\text{when }d=o(\sqrt n).\]

For the all-tournament upper bound, \paperref{sec-3}{Section 3} treats large score variance and extreme degrees separately, then retains the score penalty in the remaining class. The next lemma bounds the same spectral factor for every tournament.

\Needspace{7\baselineskip}
\subsection{A spectral cap and convex packing}\label{sec-2.4}

\paperstatement{Lemma 2.2 (Spectral cap).}{2.2}{thm-2.2} Every tournament symbol matrix satisfies
\[\|S\|_{\rm op}\le\cot\left(\frac{\pi}{2n}\right)<\frac{2n}{\pi},
\qquad \rho_n(S)\le C_*.\]

\textbf{Proof.} We include the phase-ordering proof of the operator cap \paperref{bib-4}{[4, Theorem 3.1 and Corollary 3.2]}, followed by the convex optimization needed here. The matrix \(iS\) is Hermitian with symmetric spectrum, so its largest eigenvalue is \(\|S\|_{\rm op}\). For any complex vector \(z\), apply a signed permutation simultaneously to \(z\) and \(S\) so that the nonzero coordinate phases lie in \([0,\pi)\) in increasing order. Zero coordinates are placed arbitrarily. Then \(c_{ij}=\operatorname{Im}(\overline z_i z_j)\ge0\) for \(i<j\). If \(T_n^0\) has upper-triangular entries one, then
\[z^*iSz=-2\sum_{i<j}S_{ij}c_{ij}
\le2\sum_{i<j}c_{ij}=z^*(-iT_n^0)z.\]

Taking Rayleigh maxima gives \(\|S\|_{\rm op}\le\|T_n^0\|_{\rm op}\). To compute the latter, the eigenvalue equations for adjacent rows give \((\lambda-1)v_i=(\lambda+1)v_{i+1}\). Their geometric ratio \(r=(\lambda-1)/(\lambda+1)\) satisfies \(r^n=-1\) by the first-row equation. Thus the eigenvalues are \(i\cot((2j-1)\pi/(2n))\), \(1\le j\le n\), and the claimed norm follows. For \(n\ge2\), the final strict inequality uses \(\tan u>u\); the case \(n=1\) is immediate.

Hence \(0\le x_j\le a_*\) and \(\sum_jx_j<1/2\). The function
\[\phi(x)=\log\frac{1+x}{1-x}\]

is increasing and convex on \([0,1)\). Transferring mass between two interior coordinates toward a capped coordinate and a remainder cannot decrease \(\sum_j\phi(x_j)\). Appending zeros if necessary, and increasing the total mass to \(1/2\), the relaxed maximum has coordinates \(a_*,1/2-a_*,0,\ldots\), because \(1/4<a_*<1/2\). Exponentiation gives exactly \(C_*\). The feasible spectra of tournament matrices form a subset of this relaxed region, so its maximum is an upper bound for \(\rho_n(S)\). This proves the lemma.

The same eigenvalue calculation, equivalently \paperref{bib-4}{[4, Theorem 2.1]}, also yields
\[\det(I+zT_n^0)=\frac{(1+z)^n+(1-z)^n}{2}.\]

We will use this determinant polynomial to evaluate the spectral factor of the carousel construction in \paperref{sec-4}{Section 4}.

\Needspace{9\baselineskip}
\section{The upper bound for all tournaments}\label{sec-3}

Write
\[V(T)=\sum_i\left(d_i^+-\frac{n-1}{2}\right)^2
=\frac14\|S\mathbf1\|_2^2,
\qquad
\tau=\frac{4V(T)}{(n-1)^2}.\]

We divide the proof into three cases. If \(V(T)\ge16n^2\log n\), the variance penalty in Brégman\textquotesingle s inequality makes the path count negligible. Otherwise only \(O(\log n)\) vertices can have degree outside \([(n-1)/20,19(n-1)/20]\). If any such vertices occur, we separate them from the rest and obtain a path count below the random expectation. The remaining tournaments have all normalized scores bounded away from \(\pm1\) and \(\tau=O(\log n)\). For them, the permanent estimate below retains the factor \(e^{-\tau}\) needed to obtain an \(O(n^{-1})\) final error.

Throughout this section, \(K\) may increase from one occurrence to the next. Its value is absolute once the displayed degree thresholds are fixed. The long terms have already been bounded by \(o(n^{-1})\mu_n\) in \paperref{sec-2.2}{Section 2.2}, so we estimate only the short terms \(|U|\le k_n\).

\Needspace{7\baselineskip}
\subsection{The permanent estimate used in the reduction}\label{sec-3.1}

The following lemma is the analytic input to the counting argument. Its proof, including existence of the required diagonal scaling, is given in \paperref{sec-6.4}{Section 6.4}. Row and column deletions are allowed to differ because a cycle cover may use different core vertices to enter and leave the exceptional set.

\paperstatement{Lemma 3.1 (Nonprincipal permanent bound).}{3.1}{thm-3.1} Fix \(0\le b<1\) and positive constants \(A_0,B_0\). For a tournament core of order \(N\), put \(a=S\mathbf 1/(N-1)\) and \(\tau=\sum_i a_i^2\). Suppose
\[\max_i|a_i|\le b,\qquad
\tau\le A_0\log N,\qquad
|I|=|J|=t\le B_0\log N.\]

Define \(\ell_i=(1+a_i)^{-1}\), \(r_i=(1-a_i)^{-1}\),
\[\Gamma=\prod_i(1-a_i^2),\qquad
\mathcal M_\tau=\frac{\tau+\tau^2}{N}+\frac{\tau^{3/2}}{\sqrt N}.\]

For the complementary row and column sets \(R,T\), the bound is
\[\operatorname{per}A[R,T]\le
e^{-1}D_N(S)\Gamma
\left(\prod_{i\in I}\ell_i\right)
\left(\prod_{j\in J}r_j\right)
\frac{(N-t)!}{2^{N-t}}\exp(\varepsilon_{\tau,t}),\]
\nopagebreak[3]\[\varepsilon_{\tau,t}\le K\left[
\mathcal M_\tau+\sqrt{\tau/N}+\frac{t+1}{N}
+t\sqrt{\tau/N}+\frac{t^2}{N}\right].\]

For all sufficiently large \(N\), the estimate holds simultaneously for every such tournament and every pair \(I,J\). The constant \(K\) and the threshold depend only on \(b,A_0,B_0\).

The factor \(\Gamma\) records the cost of the degree imbalance. In a principal deletion \(I=J=U\), the restored vertex weight is \(\ell_i r_i=(1-a_i^2)^{-1}\). For different deletions, the row and column factors remain separate. The error includes terms of order \(\sqrt{\tau/N}\), so retaining \(\Gamma\le e^{-\tau}\) is essential in the last case of the proof.

\Needspace{7\baselineskip}
\subsection{High score variance}\label{sec-3.2}

\paperstatement{Lemma 3.2 (Variance-sensitive permanent bound).}{3.2}{thm-3.2} There is an absolute constant \(K\) such that every tournament \(Q\) of order \(m\ge1\), with outdegrees \(d_i\) and degree variance \(V(Q)=\sum_i(d_i-(m-1)/2)^2\), satisfies
\[\operatorname{per}A_Q\le
K\sqrt{m+1}\frac{m!}{2^m}\exp\left(-\frac{V(Q)}{8m^2}\right).\]

\textbf{Proof.} We refine the Brégman bound \paperref{bib-2}{[2]} by using the concavity of its row-degree factor. If any \(d_i=0\), the asserted bound for positive order is immediate. Otherwise put \(f(k)=\log(k!)/k\) for positive integers. A direct calculation gives, for \(k\ge2\),
\[2f(k)-f(k-1)-f(k+1)
=\frac{2[\log k-f(k-1)]-k\log(1+1/k)}{k(k+1)}.\]

Arithmetic--geometric mean applied to \(1,\ldots,k-1\) gives \(\log k-f(k-1)\ge\log2\). Also \((1+1/k)^k<e<3\). The numerator is therefore at least \(\log(4/3)\ge1/4\). On the degree interval \([1,m-1]\), the piecewise-linear interpolation of \(f(k)+k^2/(8m^2)\) is concave. Jensen\textquotesingle s inequality at the mean \(\eta=(m-1)/2\) gives
\[\sum_i f(d_i)
\le m\widetilde f(\eta)-\frac{V(Q)}{8m^2}+\frac{1}{32m},\]

where \(\widetilde f\) is linear interpolation. The last term accounts for a half-integer mean and is unnecessary for an integer mean. Stirling\textquotesingle s formula gives
\[\exp(m\widetilde f(\eta))\le K\sqrt{m+1}\frac{m!}{2^m}.\]

Brégman\textquotesingle s row-degree bound now proves the variance estimate. \(\square\)

We apply this bound to each short deleted tournament. The variance must remain large after the deletion, which is why the following comparison is needed.

For a deletion set \(U\) of size \(k\), each surviving centered degree changes by at most \(k/2\). The removed squared deviations sum to at most \(kn^2/4\). Applying Cauchy--Schwarz to the cross term therefore gives
\[V(T-U)\ge V(T)-\frac{kn^2}{4}-k\sqrt{nV(T)}.\]

If \(V(T)\ge16n^2\log n\), then uniformly for \(k\le k_n\),
\[V(T-U)\ge(1-o(1))V(T).\]

The variance penalty for every short permanent is at most \(n^{-2+o(1)}\). Its factor \(\sqrt n\) is harmless; summing the determinant weights shows that the normalized short contribution is at most \(n^{-3/2+o(1)}\). Together with the long-subset tail,
\[V(T)\ge16n^2\log n\quad\Longrightarrow\quad
H(T)/\mu_n=o(1/n),\]

uniformly over this class.

\Needspace{7\baselineskip}
\subsection{Exceptional vertices}\label{sec-3.3}

Suppose henceforth that \(V(T)<16n^2\log n\). Define the exceptional set
\[F=\left\{i:\frac{d_i^+}{n-1}\notin[1/20,19/20]\right\},
\qquad f=|F|,\qquad M=T-F,\qquad N=n-f.\]

An exceptional vertex contributes at least a fixed positive multiple of \(n^2\) to \(V(T)\), so \(f=O(\log n)\). A nonexceptional full score has absolute value at most \(0.9(n-1)\). Deleting \(F\) changes it by at most \(f\), whence, uniformly for all sufficiently large \(n\),
\[\max_i|(S_M\mathbf1)_i/(N-1)|<0.95,\qquad
\tau_M=O(\log n).\]

For the squared-score bound, use \(\|S_M\mathbf1\|_2\le\|S\mathbf1\|_2+f\sqrt N\) and divide by \(N-1\). The same argument applies after any further short core deletion. Throughout those deletions we keep the original exceptional set \(F\), so the constants are uniform.

We first bound \(\operatorname{per}A_T\) by the number of ways a cycle cover can pass between \(F\) and the core \(M\). An exceptional vertex has few neighbors in one of the two directions. A cycle cover must use both directions unless it matches that vertex inside \(F\), and the latter choice will cost a factor of order \(1/N\). Use the paired core weights \(\ell,r\) of \paperref{thm-3.1}{Lemma 3.1}. Their total displacement satisfies
\[\sum_{i\in M}|\ell_i-1|+\sum_{i\in M}|r_i-1|
\le K\sqrt{N\log n}.\]

For \(x\in F\) let \(q_x=N^{-1}|\{j\in M:x\to j\}|\). Then \(q_x\) is within \(O(f/n)\) of \([0,1/20]\cup[19/20,1]\). Weighted cross-neighbor sums are bounded by \(N(q_x+\beta)\) and \(N(1-q_x+\beta)\), where \(\beta=O(\sqrt{\log n/n})\). Put
\[u_x=2(q_x+\beta),\qquad v_x=2(1-q_x+\beta).\]

Uniformly in \(x\), we can choose
\[u_xv_x\le c_n=19/100+o(1),\qquad u_x,v_x\le L_n=2+o(1).\]

Decompose a permutation contributing to the permanent by its edges inside \(F\). If there are \(s\) such edges, their row set \(R\) and column set \(C\) each have size \(s\), and each of the two cross directions has \(t=f-s\) edges. The core row and column deletion sets \(I,J\) each have size \(t\), but need not agree. The exact four-block expansion is
\[\begin{aligned}
\operatorname{per}A_T
&=\sum_{s=0}^{f}
\sum_{\substack{R,C\subseteq F\\|R|=|C|=s}}
\sum_{\substack{I,J\subseteq M\\|I|=|J|=f-s}}
\\&\quad\operatorname{per}A[R,C]\,
\operatorname{per}A[F\setminus R,J]\,
\\&\quad\operatorname{per}A[I,F\setminus C]\,
\operatorname{per}A[M\setminus I,M\setminus J].
\end{aligned}\]

Every contributing permutation is classified exactly once. In the core bound, a deleted row supplies \(\ell_i\), and a deleted column supplies \(r_j\). Absorbing these factors into the corresponding cross matchings and dropping the distinctness restrictions bounds the two cross sums by
\[N^{2t}2^{-2t}
\left(\prod_{x\notin R}u_x\right)
\left(\prod_{y\notin C}v_y\right).\]

The internal permanent is at most \(s!\). Upon division by \(n!/2^n\), the exact remaining factorial and power-of-two factor is
\[2^sN^{2t}\frac{(N-t)!}{(N+f)!}
=\left(\frac2N\right)^s\exp(O(f^2/N)).\]

Here \(N+f-(N-t)=f+t=2f-s\), so the factorial ratio supplies \(N^{-(2f-s)}\exp(O(f^2/N))\). Multiplication by \(N^{2t}\) leaves exactly \(N^{-s}\).

If \(h=|R\cap C|\), the paired product satisfies
\[\left(\prod_{x\notin R}u_x\right)
\left(\prod_{y\notin C}v_y\right)
\le c_n^{f-2s}L_n^{2s}(c_n/L_n^2)^h
\le c_n^{f-2s}L_n^{2s}.\]

The first inequality follows by counting the vertices carrying both factors, one factor, or no factor. It remains valid when \(f-2s\) is negative. The error in \paperref{thm-3.1}{Lemma 3.1} is \(o(1)\) uniformly over \(t\le f=O(\log n)\); we can discard its factor \(\Gamma\le1\). Since \(\binom fs^2s!\le f^{2s}/s!\), the complete sum gives
\[\frac{\operatorname{per}A_T}{n!/2^n}
\le(1+o(1))e^{-1}D_N(S_M)c_n^f
\exp(O(f^2/N))
\exp\left(\frac{2L_n^2f^2}{c_n^2N}\right).\]

The displayed extra exponents are \(o(1)\). All these statements remain uniform after a short core subset \(U\) is deleted. The core Gaussian factor then changes by \(\exp(O(|U|/n))\): deleting rows and columns of \(S_M/N\) costs \(O(|U|/n)\) in logarithm by the Gaussian deletion lemma, and changing normalization from \(N\) to \(N-|U|\) has the same cost.

In the path convolution, the fraction of \(k\)-subsets meeting \(F\) is at most \(fk/n\). The generic permanent bound therefore makes their entire normalized short contribution at most
\[K\frac{f}{\sqrt n}\sum_k k w_k=o(1).\]

For subsets contained in the core, restore the falling factorial and sum positive principal minors. For any core order \(N\le n\), the rank-one determinant identity gives
\[\det\left(I_N+\frac2n(I_N+A_M)\right)
\le2e\det(I_N+S_M/N).\]

To see this, extract \((1+1/n)^N\le e\). The remaining all-one rank-one factor is at most \(1+N/(n+1)<2\), since the real quadratic form of the inverse of \(I+S_M/(n+1)\) is at most the squared vector norm. Finally, the skew-frequency determinant product increases as its scale increases from \(1/(n+1)\) to \(1/N\).

For short subsets the falling-factorial restoration contributes \(\exp(O(k_n^2/n))=1+o(1)\). Because \(c_n<1/4\) eventually, the preceding uniform estimates imply
\[\frac{H(T)}{\mu_n}
\le(1/4)^f\rho_N(S_M)+o(1)
\le C_*/4+o(1)<1,\qquad f\ge1.\]

Since \(P(n)\ge\mu_n\), the fixed gap below one excludes these tournaments from the maximum for all sufficiently large \(n\). This case requires only an \(o(1)\) error.

\Needspace{7\baselineskip}
\subsection{The score penalty in the remaining class}\label{sec-3.4}

We are left with tournaments satisfying
\[\max_i|a_i|\le0.9,\qquad
a=s/(n-1),\qquad \tau=\sum_i a_i^2\le K\log n.\]

Here and below, fix one absolute constant \(K_0\) with \(\tau\le K_0\log n\). Their paired product obeys
\[\Gamma=\prod_i(1-a_i^2)\le e^{-\tau}.\]

We now use \paperref{thm-3.1}{Lemma 3.1} on the whole tournament. The score product is common to every short term; keeping it outside the sum will absorb errors that are larger than \(1/n\) individually.

For a principal deletion set \(U\) of size \(k\le k_n\), \paperref{thm-3.1}{Lemma 3.1} gives
\[\operatorname{per}A[U^c]\le
e^{-1}D_n(S)\Gamma
\left(\prod_{i\in U}g_i\right)
\frac{(n-k)!}{2^{n-k}}\exp(R_\tau+R_{\tau,k}),\]
\nopagebreak[3]\[g_i=(1-a_i^2)^{-1}\le G:=100/19,\]
\nopagebreak[3]\[R_\tau\le K[\mathcal M_\tau+\sqrt{\tau/n}+1/n],\qquad
R_{\tau,k}\le K[k\sqrt{\tau/n}+(k+k^2)/n].\]

The common factor \(\Gamma\) must be retained while errors are summed. Define nonnegative coefficients
\[c_k=(2/n)^k\sum_{|U|=k}\det(I+A[U])\prod_{i\in U}g_i.\]

They satisfy \(c_0=1\) and \(c_k\le G^kw_k\). Restoring the falling factorial contributes
\[\log\frac{n^k}{(n)_k}=O(k^2/n),\qquad k\le k_n.\]

Put \(\delta_n=\sqrt{\tau/n}+1/n\) and \(r_k=R_{\tau,k}+\log(n^k/(n)_k)\). Uniformly over \(\tau\le K_0\log n\), we have \(0\le r_k\le K\delta_n(k+k^2)\) and \(\max_{k\le k_n}r_k=o(1)\), taking the upper error budgets nonnegative. Thus
\[\sum_{k\le k_n}c_ke^{r_k}
\le\sum_{k=0}^{n}c_k+
K\delta_n\sum_{k\ge0}G^kw_k(k+k^2)
\le\left(\sum_{k=0}^{n}c_k\right)e^{K\delta_n}.\]

The last inequality uses \(\sum_k c_k\ge1\) and convergence of the indicated fixed moments. Consequently the normalized short-path sum is at most
\[\frac{e^{-1}D_n(S)\Gamma}{2}
e^{R_\tau+K\delta_n}
\det\left(I+\frac2n\operatorname{diag}(g)(I+A)\right).\]

The finite moments of \(G^kw_k\) control the summed error, rather than the largest deletion size \(k_n\). This is what prevents an extra logarithmic factor in the final bound.

We now remove the weights without losing the score penalty. Let \(W=I+(2/n)(I+A)\) and \(\Delta=(2/n)\operatorname{diag}(g-1)(I+A)\). The symmetric part of \(W\) is at least \(I\), so \(\|W^{-1}\|_{\rm op}\le1\). Since every row of \(I+A\) has Euclidean norm at most \(\sqrt n\) and \(\sum_i(g_i-1)\le K\tau\), the trace norm (the sum of singular values) can be estimated by decomposing \(\Delta\) into rank-one row matrices. Since \(\|uv^{\mathsf T}\|_*=\|u\|_2\|v\|_2\), this gives
\[\|\Delta\|_*\le K\tau/\sqrt n.\]

The determinant inequality \(|\det(I+E)|\le\exp(\|E\|_*)\) therefore implies
\[\det(W+\Delta)\le\det(W)e^{K\tau/\sqrt n}.\]

Both determinants are positive: \(W\) has positive-definite symmetric part, and the weighted determinant has a principal-minor expansion with positive coefficients. The unweighted rank-one bound from \paperref{sec-3.3}{Section 3.3}, now with \(N=n\), is \(\det(W)\le2e\det(I+S/n)\).

Combining these inequalities and adding only the long-subset tail gives the uniform score-sensitive bound
\[\frac{H(T)}{\mu_n}\le
\rho_n(S)\exp\left\{-\tau+
K\left[\frac{1+\tau+\tau^2}{n}
+\frac{\tau^{3/2}+\tau}{\sqrt n}+\sqrt{\tau/n}\right]\right\}
+o(1/n).\]

For \(\tau\le K_0\log n\), all polynomial terms except the constant \(1/n\) can consume at most \(\tau/4\) for sufficiently large \(n\). The remaining square-root term is bounded by Young\textquotesingle s inequality:
\[K\sqrt{\tau/n}\le\tau/4+K^2/n.\]

The entire exponent is consequently at most \(-\tau/2+K/n\). The spectral cap now gives
\[H(T)/\mu_n\le C_*e^{K/n}+o(1/n)=C_*+O(1/n).\]

\Needspace{7\baselineskip}
\subsection{Completion of the upper bound}\label{sec-3.5}

The three classes cover every tournament. High-variance tournaments have normalized path count \(o(1/n)\); the low-variance tournaments with exceptional vertices have count strictly below one; all remaining tournaments satisfy the preceding \(C_*+O(1/n)\) bound. It follows that
\[\boxed{P(n)\le\bigl(C_*+O(1/n)\bigr)\mu_n.}\]

The thresholds in the three cases are independent of the tournament, so one common threshold gives the upper half of \paperref{thm-1.1}{Theorem 1.1}.

\Needspace{9\baselineskip}
\section{Small-score paths and the lower bound}\label{sec-4}

\paperstatement{Theorem 4.1 (Small-score path approximation).}{4.1}{thm-4.1} There is an absolute constant \(C\) with the following property. For every nonnegative function \(d=d(n)\) satisfying \(d(n)/\sqrt n\to0\), there is a threshold \(n_0(d)\) such that every tournament of order \(n\ge n_0(d)\) with \(\|S\mathbf1\|_\infty\le d(n)\) satisfies
\[\left|\frac{H(T)}{\mu_n}-\rho_n(S)\right|\le C\frac{(d(n)+1)^2}{n}.\]

\Needspace{7\baselineskip}
\subsection{Proof of the small-score approximation}\label{sec-4.1}

\textbf{Proof.} Write \(d=d(n)\). Use \paperref{thm-6.4}{Lemma 6.4} for every principal deletion \(|U|=k\le k_n\). Since \(d+k_n+1=o(\sqrt n)\), it applies uniformly. Restoring the falling factorial \((n)_k\), the positive path convolution gives
\[\frac{H(T)}{\mu_n}
=\frac{e^{-1}D_n(S)}2
\sum_{k\le k_n}\left(\frac2n\right)^k
\sum_{|U|=k}\det(I+A[U])
+O\left(\frac{(d+1)^2}{n}\right).\]

The error has this rate, without an additional logarithmic factor, because it is bounded by a fixed constant times
\[\frac1n\sum_{k\ge0}w_k(d+k+1)^2
=O\left(\frac{(d+1)^2}{n}\right).\]

The falling-factorial error contributes \(O(n^{-1}\sum k^2w_k)\), and the terms indexed by large deletion sets are \(o(n^{-1})\) by \paperref{sec-2.2}{Section 2.2}. Extending the displayed short generating sum to all \(k\) costs \(n^{-2+o(1)}\), since its coefficients are bounded by \(w_k\). Thus
\[\frac{H(T)}{\mu_n}
=\frac{e^{-1}D_n(S)}2
\det\left(I+\frac2n(I+A)\right)
+O\left(\frac{(d+1)^2}{n}\right).\]

Put \(u=\mathbf1/\sqrt n\) and \(Q=S/(n+1)\). The rank-one determinant lemma yields
\[\det\left(I+\frac2n(I+A)\right)
=(1+1/n)^n\det(I+Q)
\left(1+\frac n{n+1}u^{\mathsf T}(I+Q)^{-1}u\right).\]

Skew symmetry implies
\[u^{\mathsf T}(I+Q)^{-1}u=u^{\mathsf T}(I-Q^2)^{-1}u,\]

and its difference from one in absolute value is at most
\[\|Qu\|_2^2\le\frac{d^2}{(n+1)^2}.\]

Also \((1+1/n)^n=e(1+O(n^{-1}))\), while replacing \(S/(n+1)\) by \(S/n\) changes the log determinant by \(O(n^{-1})\), using the bounded sum of normalized squared frequencies. Consequently,
\[\det\left(I+\frac2n(I+A)\right)
=2e\det(I+S/n)\left(1+O(n^{-1}+d^2/n^2)\right).\]

Substitution proves the theorem. The estimates use one absolute constant; only the order after which the small-score conditions hold depends on the function \(d\). \(\square\)

\Needspace{7\baselineskip}
\subsection{Carousel tournaments}\label{sec-4.2}

\paperstatement{Corollary 4.2 (Carousel lower bound).}{4.2}{thm-4.2} For the odd-order carousel and the even-order one-vertex deletion defined below,
\[H(\mathrm{Car}_n)/\mu_n=L+O(n^{-1}).\]

\textbf{Proof.} For odd \(n\), define the carousel tournament \(\mathrm{Car}_n\) on residues modulo \(n\) by \(i\to i+j\) for \(1\le j\le(n-1)/2\). It is regular, so \(d=0\). For even \(n\), let \(\mathrm{Car}_n\) be a one-vertex deletion from \(\mathrm{Car}_{n+1}\); then \(d=1\).

The odd carousel symbol matrix is signed-permutation similar to \(T_n^0\). One explicit construction is to conjugate \(T_n^0\) by the diagonal signs \((-1)^i\), \(0\le i<n\), and reorder its indices as \(0,2,\ldots,n-1,1,3,\ldots,n-2\). For two indices of the same parity the switched edge follows their increasing order; for two indices of opposite parity it follows the reverse order. Reading the indices in the displayed cyclic order gives precisely the carousel orientation. Restrict the inverse signed conjugation and reordering to an even-order one-vertex deletion: the resulting transformed matrix is \(T_n^0\) on the remaining ordered indices. Thus the same signed-spectral computation applies after deleting any one vertex. Hence both parities have
\[\rho_n(S_{\mathrm{Car}_n})
=r_n:=\frac{(n+1)^n+(n-1)^n}{(n+i)^n+(n-i)^n}.\]

Dividing numerator and denominator by \(n^n\), Taylor expansion of \((1+z/n)^n\) at the four fixed values \(z=1,-1,i,-i\) gives
\[r_n=\frac{\cosh1}{\cos1}+O(n^{-1})=L+O(n^{-1}).\]

Therefore \(H(\mathrm{Car}_n)/\mu_n=L+O(n^{-1})\) for both parities, and
\[P(n)\ge(L-O(n^{-1}))\mu_n.\]

The passage from the signed similarity to the path count uses \paperref{thm-4.1}{Theorem 4.1}, applied to scores \(d=0\) and \(d=1\). Let \(K_l,N_l\) and \(K_u,N_u\) be the constants and thresholds from the lower and upper bounds. Taking \(K=K_l+K_u\) and \(n_0=\max(N_l,N_u,2)\) proves \paperref{thm-1.1}{Theorem 1.1}. \(\square\)

\Needspace{9\baselineskip}
\section{A uniform permanent approximation}\label{sec-5}

We now prove the permanent estimate underlying the reduction. After a matrix has been scaled to have all row and column sums equal to one, it has the form \((J_n+E)/n\) with \(E\) centered on both sides. The theorem below approximates its permanent by a Gaussian determinant. It uses a bound on the entries and a fixed singular-value gap, which \paperref{sec-6}{Section 6} verifies for the scaled tournament submatrices.

McCullagh \paperref{bib-5}{[5]} obtained the determinantal leading term under moderate-deviation conditions. We give a self-contained version with fixed entry bound \(C\), singular-value gap \(q\), and an error uniform over the matrix class for each fixed \(C,q\). The coefficient and tail estimates are kept separate so that the conclusion can be applied uniformly to the deletions and scalings in \paperref{sec-6}{Section 6}. A related recent approximation by Li \paperref{bib-8}{[8, Theorem 2.1]} assumes that the maximum absolute row or column sum of the centered perturbation is \(o(n)\); that hypothesis does not cover dense kernels with such sums of order \(n\).

Write \(P_n=J_n/n\). The operator norm is the Euclidean operator norm. For a complex matrix \(Z\), write \(Z^*\) for its conjugate transpose and \(|Z|=(Z^*Z)^{1/2}\) for its operator absolute value.

\Needspace{7\baselineskip}
\subsection{Statement and proof plan}\label{sec-5.1}

\paperstatement{Theorem 5.1 (Uniform permanent approximation).}{5.1}{thm-5.1} Fix \(0\le C<\infty\) and \(0\le q<1\). For \(n\ge1\), suppose that \(E\in\mathbb R^{n\times n}\) satisfies
\[E\mathbf1=E^{\mathsf T}\mathbf1=0,\qquad
\max_{i,j}|E_{ij}|\le C,\qquad
\|E/n\|_{\mathrm{op}}\le q.\]

Set \(B=E/n\). Uniformly over all such matrices,
\[\frac{\operatorname{per}(J_n+E)}{n!}
=\det(I-BB^{\mathsf T})^{-1/2}+O_{C,q}(n^{-1}).\]

The same statement holds with a relative factor \(1+O_{C,q}(n^{-1})\) multiplying the determinant factor.

Here and below the determinant square root is the positive square root on the real interval under consideration. The assumptions imply \(\|B\|_{\mathrm F}^2\le C^2\). Consequently, the determinant factor is bounded above by a constant depending only on \(C,q\), and it is at least one.

The leading term comes from pairings. Inclusion--exclusion on repeated row and column indices turns each permanent coefficient into a sum of bipartite multigraphs. Centering removes degree-one vertices, and the degree-two components sum to the Gaussian determinant. The remaining components are smaller by powers of \(n^{-1}\), measured by their edge excess. This proves the approximation through a linear range of degrees. A separate complex-analytic bound then controls the rest of the actual permanent polynomial.

For the proof, set
\[f(t)=\frac{\operatorname{per}(J_n+tE)}{n!},
\qquad G(t)=\det(I-t^2BB^{\mathsf T})^{-1/2}.\]

Choose \(1<\sigma<R<1/q\), with no upper restriction when \(q=0\); for example,
\[R=\frac{3+q}{2(1+q)},\qquad \sigma=\frac{1+R}{2}.\]

We use the envelope
\[\Gamma(r)=\exp\left(\frac{C^2r^2}{2(1-q^2r^2)}\right),
\qquad r\ge0,\ rq<1.\]

Throughout this section, upper restrictions involving \(1/q\) are omitted when \(q=0\). All constants depend only on \(C,q\) and the chosen radii. \paperref{sec-5.6}{Section 5.6} collects them into an explicit error bound.

\Needspace{7\baselineskip}
\subsection{Coefficients and Gaussian pairings}\label{sec-5.2}

Write \(f(t)=\sum_{k=0}^n a_kt^k\). Expansion of the permanent by the entries from \(E\) gives
\[a_k=\frac1{(n)_k}
\sum_{\substack{I,J\subseteq[n]\\|I|=|J|=k}}
\operatorname{per}E[I,J],\]

where \((n)_k=n(n-1)\cdots(n-k+1)\) and \((n)_0=1\). Define
\[\begin{gathered}
F_k=\frac{(n)_k}{n^k}a_k
=\frac1{k!}
\sum_{\substack{i_1,\ldots,i_k\ \mathrm{distinct}\\
j_1,\ldots,j_k\ \mathrm{distinct}}}
\prod_{\ell=1}^kB_{i_\ell j_\ell},
\\
F(t)=\sum_{k=0}^nF_kt^k.
\end{gathered}\]

The right-hand sum also makes sense for \(k>n\), when it is empty and equals zero. This convention is useful for the formal identities below.

Let \(\Pi_k\) be the lattice of set partitions of \([k]\). The indicator that \(k\) coordinates are distinct has the standard partition-lattice inclusion--exclusion expansion. Paired partition expansions also underlie the permanent calculations in \paperref{bib-5}{[5, Section 4]}. The Möbius weight of a block of size \(d\) is \((-1)^{d-1}(d-1)!\). Apply this expansion independently to the row and column coordinates in \(F_k\). A pair of partitions becomes a bipartite multigraph: its edges carry labels \(1,\ldots,k\), its row and column vertices are the partition blocks, and each block of size \(d\) has the above weight.

The numerical label of each vertex is summed independently over \([n]\); numerical labels of different vertices are allowed to coincide. If a row vertex has degree one, summing its label gives a column sum of \(B\), which is zero. A degree-one column vertex similarly gives a zero row sum. Hence only graphs with all vertex degrees at least two survive.

For example, in degrees two and three the only surviving row and column partitions each consist of one block. Their paired Möbius weights are \(1\) and \(4\), respectively; dividing by \(k!\) gives
\[F_2=\frac12\sum_{i,j}B_{ij}^2,
\qquad
F_3=\frac23\sum_{i,j}B_{ij}^3,
\qquad |F_3|\le\frac{2C^3}{3n}.\]

The quadratic term belongs to the Gaussian factor below. The cubic term is the first possible core correction, and its \(n^{-1}\) bound illustrates the excess estimate.

The connected components in which every vertex has degree two are the pure two-degree components. Their total exponential generating function is \(G(t)\). One direct verification uses two independent standard real Gaussian vectors \(X,Y\). Wick\textquotesingle s formula shows that the coefficient of \(t^k\) in
\[\mathbb E\exp(tX^{\mathsf T}BY)\]

is exactly the sum over the two pairing partitions on the edge labels. These are the graphs all of whose vertices have degree two. Conditional Gaussian integration, followed by diagonalization of \(BB^{\mathsf T}\), gives
\[\mathbb E\exp(tX^{\mathsf T}BY)
=\mathbb E\exp\!\left(\frac{t^2}{2}X^{\mathsf T}BB^{\mathsf T}X\right)
=\det(I-t^2BB^{\mathsf T})^{-1/2}.\]

The analytic identity holds for \(|t|\|B\|_{\mathrm{op}}<1\) and identifies all its formal coefficients. In particular, \(G\) has nonnegative coefficients, and
\[G(r)\le\Gamma(r)\qquad(r\ge0,\ rq<1),\]

because \(-\log(1-x)\le x/(1-x)\) and \(\sum_i s_i^2=\|B\|_{\mathrm F}^2\le C^2\), where \(s_i\) are the singular values of \(B\).

Remove every pure two-degree component from a surviving graph. Call what remains its core, which may be disconnected. If the core has \(h\) edges and \(v\) vertices, define its excess by \(j=h-v\). A nonempty core has \(j\ge1\). The labeled component decomposition gives the coefficientwise identity
\[F(t)=G(t)\left(1+\sum_{j\ge1}C_j(t)\right),\]

where \(C_j\) is the generating series of cores of excess \(j\), retaining all Möbius signs and matrix contractions. This is a formal identity. In a fixed degree \(k\), only finitely many excesses occur; in fact, a nonempty core contributing to degree \(k\) has \(j<k\). We never assume that the infinite sum over \(j\) converges at a nonzero value of \(t\). Above degree \(n\), its coefficients cancel to zero because the distinct-coordinate definition of \(F_k\) is zero.

\Needspace{7\baselineskip}
\subsection{Bounding the non-Gaussian cores}\label{sec-5.3}

A long chain of degree-two vertices can be summed using the operator gap. Compressing these chains leaves at most \(2j\) vertices at excess \(j\), so the remaining counting problem depends on \(j\) rather than on the original degree. Set
\[W=\max\left\{1,CR+\frac{C^2R^2}{1-Rq}\right\},
\qquad D=710W^3,\qquad
T_1=\frac32W^2+\frac{10}{3}W^3.\]

\paperstatement{Lemma 5.2 (Compressed-core activity).}{5.2}{thm-5.2} With the notation above,
\[\|C_j\|_R\le(Dj/n)^j\qquad(j\ge1),\]

and the first excess satisfies the sharper bound
\[\|C_1\|_R\le T_1/n.\]

Here \(\|H\|_R=\sum_{k\ge0}|[t^k]H|R^k\).

\emph{Proof.} Compress every maximal path whose internal vertices have degree two. Each resulting edge is a positive-length chain between vertices of degree at least three. Loops, parallel edges, and disconnected compressed graphs are allowed. If \(b\) is the number of high-degree vertices and \(e\) the number of compressed edges, compression preserves excess and yields
\[e=b+j,\qquad
2j=\sum_{i=1}^b(d_i-2),\qquad
1\le b\le2j,\qquad e\le3j.\]

The following decorated count determines the symmetry factors. Temporarily label the \(b\) high-degree vertices and the \(d_i\) half-edges at each vertex. Use these labels to select one reading direction for every chain. A given original edge-labeled core has \(b!\prod_i d_i!\) such decorations. Conversely, a decorated half-edge pairing and a list of positive chain lengths determine all chain positions. The original \(h\) edge labels can be assigned to these positions in \(h!\) ways. Subject to the bipartite parity condition, the internal vertices and their colors are then determined.

After division by the original exponential generating-function factor \(h!\) and removal of the decorations, multiplication by the absolute Möbius weights \(\prod_i(d_i-1)!\) leaves
\[\frac1{b!\prod_i d_i}.\]

This argument includes loops and chain reversals: their two half-edges are already labeled, so no extra direction factor is introduced. Internal degree-two vertices have absolute Möbius weight one. There is no additional chain-length factorial.

Ignore the color restrictions only when taking an upper bound. For fixed \(j,b\), the resulting total compensation is
\[U_{j,b}=
\frac{2^b}{b!}\frac{(2e)!}{2^e e!}
\sum_{\substack{d_1+\cdots+d_b=2e\\d_i\ge3}}
\frac1{d_1\cdots d_b},
\qquad e=b+j.\]

For a chain of length one the matrix entry is bounded by \(C/n\). For a chain of length \(\ell\ge2\), first sum its internal labels. Its contraction is an entry of an alternating product of \(B\) and \(B^{\mathsf T}\). Each endpoint row or column has norm at most \(C/\sqrt n\), and all intermediate factors have operator norm at most \(q\). Thus the absolute contraction is at most
\[\frac{C^2}{n}q^{\ell-2}.\]

Summing over lengths with \(R\) weights bounds one chain by \(W/n\). Summing the \(b\) high-degree labels contributes at most \(n^b\). Hence
\[\|C_j\|_R\le n^{-j}
\sum_{b=1}^{2j}U_{j,b}W^{b+j}.\]

The number of degree sequences is \(\binom{2j-1}{b-1}\le4^j\): write \(d_i=3+x_i\) and \(\sum_i x_i=2j-b\). Also \(\prod_i d_i^{-1}\le3^{-b}\), and the number of pairings is at most \((6j)^{b+j}\). It follows that
\[\sum_{b=1}^{2j}U_{j,b}W^{b+j}
\le(24jW)^j
\sum_{b=1}^{2j}\frac{(4jW)^b}{b!}.\]

For \(\theta=1/(2W)\le1/2\),
\[\sum_{b=0}^{2j}\frac{(4jW)^b}{b!}
\le\theta^{-2j}\exp(4jW\theta)
=(2\mathrm e W)^{2j}.\]

Therefore the activity is at most \((96\mathrm e^2W^3j)^j\), which is bounded by \((710W^3j)^j\). For completeness, \(\mathrm e<87/32\) follows by summing its exponential series through degree six and bounding the tail by \(8/(7\cdot7!)\); this rational bound gives \(96\mathrm e^2<710\).

For \(j=1\) only \(b=1,2\) are possible. The compensation formula gives
\[U_{1,1}=\frac32,\qquad U_{1,2}=\frac{10}{3}.\]

Their chain bounds give \(\|C_1\|_R\le T_1/n\). \(\square\)

\Needspace{7\baselineskip}
\subsection{A linear coefficient window and factorial recovery}\label{sec-5.4}

The bound in \paperref{thm-5.2}{Lemma 5.2} is useful while \(j/n\) is small. Choose
\[\alpha=\min\left\{\frac14,\frac1{16D},\frac{\log\sigma}{2}\right\},
\qquad M=\lfloor\alpha n\rfloor.\]

Put \(v_j=(Dj/n)^j\). Whenever \(j+1\le M\),
\[\frac{v_{j+1}}{v_j}
=\frac{D(j+1)}n\left(1+\frac1j\right)^j
\le\frac{\mathrm e D(j+1)}n
\le3D\alpha<\frac12.\]

Thus, for \(n\ge4/\alpha\),
\[\sum_{j=2}^{M}v_j\le2v_2=\frac{8D^2}{n^2}.\]

Only excesses \(j\le M\) can contribute to degrees at most \(M\) in the formal core identity. Positivity of the coefficients of \(G\), the activity lemma, and the convolution inequality for weighted coefficient sums therefore give
\[\sum_{k=0}^{M}|F_k-[t^k]G|R^k
\le\Gamma(R)\left(\frac{T_1}n+\frac{8D^2}{n^2}\right).\]

This bound concerns a finite coefficient window, not the sum of all excesses evaluated at \(R\).

Let \(g_k=[t^k]G\) and \(r_k=n^k/(n)_k\) for \(0\le k\le M\). Since \(\alpha\le1/4\),
\[\log r_k
=\sum_{\ell=0}^{k-1}-\log(1-\ell/n)
\le\frac{k(k-1)}{2n(1-\alpha)}
\le\frac{k^2}n
\le\alpha k.\]

The choice \(\alpha\le(\log\sigma)/2\) implies \(r_k\le\sigma^k\), and
\[0\le r_k-1
\le\frac{k(k-1)}{2n(1-\alpha)}r_k.\]

Because \(a_k=r_kF_k\), we can split
\[a_k-g_k=r_k(F_k-g_k)+(r_k-1)g_k.\]

The sum of the first terms is bounded by the preceding \(R\)-weighted estimate. For the second terms, use the exact Gaussian coefficient moment rather than a coarse coefficient bound:
\[\sum_{k\ge0}k(k-1)g_k\sigma^k
=\left((t\partial_t)^2-t\partial_t\right)G(t)\big|_{t=\sigma}.\]

Set \(a_i=\sigma^2s_i^2\). Differentiation of \(G(t)=\prod_i(1-t^2s_i^2)^{-1/2}\) yields
\[\frac{\left((t\partial_t)^2-t\partial_t\right)G(t)}{G(t)}
\bigg|_{t=\sigma}
=\left(\sum_i\frac{a_i}{1-a_i}\right)^2
+\sum_i\frac{a_i(1+a_i)}{(1-a_i)^2}.\]

Define
\[u_\sigma=\frac{\sigma^2C^2}{1-\sigma^2q^2},
\qquad
v_\sigma=\frac{\sigma^2C^2(1+\sigma^2q^2)}
{(1-\sigma^2q^2)^2},\]
\nopagebreak[3]\[K_G=\frac{\Gamma(\sigma)}{2(1-\alpha)}
\left(u_\sigma^2+v_\sigma\right).\]

The two terms in the derivative formula are bounded by \(u_\sigma^2\) and \(v_\sigma\), respectively. Consequently,
\[\sum_{k=0}^{M}(r_k-1)g_k\le K_G/n.\]

Combining the two recovered contributions gives the first two terms of the explicit error budget.

\Needspace{7\baselineskip}
\subsection{Controlling the remaining coefficients}\label{sec-5.5}

We have controlled the coefficients through degree \(M\), where \(M\) is proportional to \(n\). To estimate the higher coefficients by Cauchy\textquotesingle s inequality, it is enough to bound \(f\) on the fixed circle \(|t|=R>1\) by \(\exp(O(\sqrt n))\). The next two lemmas give that bound with the normalization \(n!/n^n\) intact.

\paperstatement{Lemma 5.3 (Permanent polarization).}{5.3}{thm-5.3} For every complex square matrix \(Z\),
\[|\operatorname{per}Z|
\le\sqrt{\operatorname{per}|Z|\,
\operatorname{per}|Z^*|}.\]

\emph{Proof.} In \((\mathbb C^n)^{\otimes n}\), let
\[\xi=(n!)^{-1/2}
\sum_{\pi\in S_n}
e_{\pi(1)}\otimes\cdots\otimes e_{\pi(n)}.\]

Then \(\|\xi\|=1\) and \(\langle\xi,Z^{\otimes n}\xi\rangle=\operatorname{per}Z\). Extend the polar factor in \(Z=U|Z|\) to a unitary matrix \(U\) if \(Z\) is singular. Put \(T=|Z|^{\otimes n}\) and \(V=U^{\otimes n}\). Cauchy--Schwarz in the positive semidefinite form induced by \(T\) gives
\[|\langle\xi,VT\xi\rangle|
\le\sqrt{\langle V^*\xi,TV^*\xi\rangle
\langle\xi,T\xi\rangle}.\]

The two factors are \(\operatorname{per}(U|Z|U^*)=\operatorname{per}|Z^*|\) and \(\operatorname{per}|Z|\). Both are nonnegative because they are diagonal matrix elements of positive semidefinite tensor powers. \(\square\)

\paperstatement{Lemma 5.4 (Positive semidefinite permanent bound).}{5.4}{thm-5.4} If \(H\) is Hermitian positive semidefinite with eigenvalues \(\lambda_1,\ldots,\lambda_n\), then
\[\operatorname{per}H\le\frac{n!}{n^n}h_n(\lambda_1,\ldots,\lambda_n),\]

where \(h_n\) is the complete homogeneous symmetric polynomial of degree \(n\). If \(\lambda_1=1\) and \(0\le\lambda_i<1\) for \(i\ge2\), then
\[\operatorname{per}H\le\frac{n!}{n^n}
\prod_{i=2}^n(1-\lambda_i)^{-1}.\]

\emph{Proof.} Let \(Z\) be a circular complex Gaussian vector with covariance \(H\). Complex Wick\textquotesingle s formula gives \(\operatorname{per}H=\mathbb E\prod_i|Z_i|^2\). Pointwise AM--GM bounds the product by \(n^{-n}\|Z\|^{2n}\). After unitary diagonalization, \(\|Z\|^2\) has the distribution of \(\sum_i\lambda_i|\zeta_i|^2\), where the \(\zeta_i\) are independent standard circular complex Gaussian variables. Their moments satisfy \(\mathbb E|\zeta_i|^{2k}=k!\). The multinomial expansion consequently gives
\[\mathbb E\left(\sum_i\lambda_i|\zeta_i|^2\right)^n
=n!h_n(\lambda_1,\ldots,\lambda_n).\]

Finally, \(h_n(1,\lambda_2,\ldots,\lambda_n)\) is the sum of all monomials in \(\lambda_2,\ldots,\lambda_n\) of total degree at most \(n\), and is at most the product of their infinite geometric sums. \(\square\)

The Gaussian identity and AM--GM comparison in \paperref{thm-5.4}{Lemma 5.4} are the tools used in Han--Niles-Weed \paperref{bib-3}{[3, Lemmas 4.3–4.4]} for positive semidefinite inputs. \paperref{thm-5.3}{Lemma 5.3} connects them to the general matrix required here.

Put \(\beta=RC/(1-Rq)\). The centering assumptions give \(P_nB=BP_n=0\). On \(|t|=R\) they imply
\[|P_n+tB|=P_n+R(B^{\mathsf T}B)^{1/2},
\qquad
|(P_n+tB)^*|=P_n+R(BB^{\mathsf T})^{1/2}.\]

Both matrices have eigenvalues \(1\) and \(Rs_i\) on \(\mathbf1^\perp\), including zero singular values as needed. Since \(Rq<1\), the two lemmas and the exact identity \(f(t)=(n^n/n!)\operatorname{per}(P_n+tB)\) give
\[|f(t)|\le\prod_i(1-Rs_i)^{-1}
\le\exp\!\left(\frac{R}{1-Rq}\sum_i s_i\right)
\le\exp(\beta\sqrt n).\]

The last step uses \(\sum_i s_i\le\sqrt n\,\|B\|_{\mathrm F}\le C\sqrt n\). Thus Cauchy\textquotesingle s coefficient estimate gives
\[\sum_{k>M}|a_k|
\le\frac{R}{R-1}
\exp\!\left(\beta\sqrt n-\alpha n\log R\right).\]

For the Gaussian tail, positivity of its coefficients gives the convenient, slightly weaker bound
\[\sum_{k>M}g_k
\le \Gamma(\sigma)\sigma^{-M}
\le\sigma\Gamma(\sigma)\exp\!\left(-\alpha n\log\sigma\right).\]

The two tails are exponentially small in \(n\). Combined with the \(O(n^{-1})\) finite-window estimate, they prove \paperref{thm-5.1}{Theorem 5.1}. \(\square\)

\Needspace{7\baselineskip}
\subsection{The combined error bound}\label{sec-5.6}

\paperstatement{Proposition 5.5 (Explicit error bound).}{5.5}{thm-5.5} For every integer \(n\ge4/\alpha\),
\[\begin{aligned}
|f(1)-G(1)|
&\le \frac{\Gamma(R)T_1+K_G}{n}
+\frac{8\Gamma(R)D^2}{n^2}
\\&\quad+\frac{R}{R-1}\exp\!\left(\beta\sqrt n-\alpha n\log R\right)
+\sigma\Gamma(\sigma)\exp\!\left(-\alpha n\log\sigma\right).
\end{aligned}\]

The first two terms are the coefficient-window and factorial-recovery errors from \paperref{sec-5.4}{Section 5.4}; the last two are the tails from \paperref{sec-5.5}{Section 5.5}. The exponential terms are \(o(n^{-1})\), proving the uniform rate in \paperref{thm-5.1}{Theorem 5.1}.

\Needspace{9\baselineskip}
\section{Scaling and nonprincipal permanent estimates}\label{sec-6}

\paperref{thm-5.1}{Theorem 5.1} applies to a matrix with equal row and column sums after a scalar normalization. A tournament matrix generally has unequal margins: for \(C=2A/(n-1)\),
\[C\mathbf1=\mathbf1+a,\qquad C^{\mathsf T}\mathbf1=\mathbf1-a.\]

We must correct these margins while controlling the change in the permanent and in its Gaussian factor. We first construct a local diagonal scaling and bound its total logarithmic cost. A Gram-matrix comparison then controls the effect of deleting rows and columns. Finally, paired score corrections put the tournament submatrices in the domain of the local theorem and yield \paperref{thm-3.1}{Lemma 3.1}.

Write \(\mathbf1_p\) for the all-ones vector, \(P_p=\mathbf1_p\mathbf1_p^{\mathsf T}/p\), and \(\Pi_p=I-P_p\). The total mass of a matrix is \(\mathfrak m(X)=\sum_{i,j}X_{ij}\). We use \(\|\cdot\|_{\rm op}\) for the Euclidean operator norm, \(\|\cdot\|_F\) for the Frobenius norm, and \(\|\cdot\|_*\) for the trace norm. Constants in this section are uniform once the stated density and gap parameters are fixed.

\Needspace{7\baselineskip}
\subsection{Local scaling and its cost}\label{sec-6.1}

We seek \(B_{ij}=X_{ij}e^{x_i+y_j}\) with all margins equal to one. Replacing \((x,y)\) by \((x+c\mathbf1,y-c\mathbf1)\) leaves \(B\) unchanged; the condition \(\sum x_i=\sum y_j\) removes this one-dimensional freedom. The singular-value gap controls the linearized balancing equations on the remaining subspace. The entry bound upgrades that control to an infinity-norm estimate independent of the dimension, allowing a contraction argument.

\paperstatement{Lemma 6.1 (Local scaling).}{6.1}{thm-6.1} Let \(p\ge1\) and let \(X\) be a real \(p\times p\) matrix of mass \(p\). Put
\[\alpha=X\mathbf 1_p-\mathbf 1_p,\qquad \beta=X^{\mathsf T}\mathbf 1_p-\mathbf 1_p,\qquad \varepsilon=\max(\|\alpha\|_\infty,\|\beta\|_\infty).\]

Assume that \(|X_{ij}|\le C/p\), that each absolute row and column sum is at most \(2\), and that
\[X_0=X-\frac{\alpha\mathbf 1_p^{\mathsf T}+\mathbf 1_p\beta^{\mathsf T}}p,\qquad E=X_0-P_p\]

satisfies \(|E_{ij}|\le C_0/p\) and \(\|E\|_{\rm op}\le q<1\), where \(C,C_0,q\) are fixed. Define
\[L_0=\frac32+C_0+\frac{C_0^2}{1-q},\]
\nopagebreak[3]\[\varepsilon_0=\min\left\{\frac1{256L_0^2},\frac{1-q}{2(32L_0+2)}\right\}.\]

If \(\varepsilon\le\varepsilon_0\), there are finite real vectors \(x,y\) with \(\sum_i x_i=\sum_j y_j\) such that \(B=\operatorname{diag}(e^x)X\operatorname{diag}(e^y)\) has all row and column sums equal to \(1\). Moreover,
\[\|(x,y)\|_\infty\le4L_0\varepsilon,\qquad |B_{ij}|\le2C/p,\qquad \|B-P_p\|_{\rm op}\le(1+q)/2.\]

If \(X\ge0\), the matrix \(B\) is a genuine doubly stochastic scaling and has the same zero support as \(X\).

\textbf{Proof.} The two marginal error vectors have zero sum, so \(X_0\) has unit row and column sums, and \(E\) is centered on both sides. Let \(u=(\mathbf 1_p,-\mathbf 1_p)\) and \(N=uu^{\mathsf T}/(2p)\). The balanced Hessian \(\mathsf H_0\) has diagonal blocks \(I\) and off-diagonal blocks \(X_0,X_0^{\mathsf T}\). Its gauge-fixed inverse can be written explicitly. Set
\[R_L=(I-EE^{\mathsf T})^{-1},\qquad R_R=(I-E^{\mathsf T}E)^{-1},\qquad O=ER_R.\]

The four blocks of \((\mathsf H_0+N)^{-1}\) are
\[M_{LL}=R_L-P_p/4,\qquad M_{RR}=R_R-P_p/4,\]
\nopagebreak[3]\[M_{LR}=-O-P_p/4,\qquad M_{RL}=-O^{\mathsf T}-P_p/4.\]

Indeed, on the common constant direction the inverse eigenvalue is \(1/2\), on the gauge direction it is \(1\), and on the two zero-sum spaces this is the standard block inverse. The dense entry bound gives
\[|(R_L-I)_{ij}|,\ |(R_R-I)_{ij}|\le\frac{C_0^2}{p(1-q^2)},\]
\nopagebreak[3]\[|O_{ij}|\le\frac{C_0}{p}+\frac{C_0^2q}{p(1-q^2)}.\]

For the last bound, separate \(O=E+EE^{\mathsf T}ER_R\) and use the Euclidean norms of the endpoint row and column. Consequently both induced one- and infinity-norms of this inverse are at most \(L_0\).

Let \(\mathsf H_X\) have diagonal blocks \(\operatorname{diag}(X\mathbf 1_p)\) and \(\operatorname{diag}(X^{\mathsf T}\mathbf 1_p)\), and off-diagonal blocks \(X,X^{\mathsf T}\). The difference \(\mathsf H_X-\mathsf H_0\) has both induced norms at most \(3\varepsilon\). A Neumann series therefore bounds both induced norms of \((\mathsf H_X+N)^{-1}\) by \(2L_0\).

Writing \(z=(x,y)\) and \(g=(\alpha,\beta)\), the balance equations are
\[0=g+\mathsf H_Xz+\mathcal N_X(z).\]

Here \(\mathcal N_X(z)\) is the concatenation of the row and column sums of \(X_{ij}(e^{x_i+y_j}-1-x_i-y_j)\). On a ball \(\|z\|_\infty\le R\le1/4\), the absolute margin assumptions imply
\[\|\mathcal N_X(z)\|_\infty\le8R^2,\qquad \|D\mathcal N_X(z)\|_{\infty\to\infty}\le16R.\]

All these vectors are perpendicular to \(u\). On the gauge subspace \(u^\perp\), consider
\[\mathcal T(z)=-(\mathsf H_X+N)^{-1}(g+\mathcal N_X(z)).\]

Take \(R=4L_0\varepsilon\). The image of the origin has norm at most \(R/2\), and the Lipschitz constant is at most \(128L_0^2\varepsilon\le1/2\). The map preserves the gauge ball and is contractive, giving the asserted finite potentials. If \(\varepsilon=0\), take \(z=0\) directly.

Finally, \(R\le1/4\) implies the entry bound. The row and column absolute sums give
\[\|B-X\|_{1\to1},\ \|B-X\|_{\infty\to\infty}\le32L_0\varepsilon.\]

Since \(\|X-X_0\|_{\rm op}\le2\varepsilon\), the second condition in \(\varepsilon_0\) proves the centered gap. Positive diagonal factors retain both nonnegativity and the zero support. This proves the lemma.

The balance point alone is insufficient for the permanent estimate: restoring the diagonal factors multiplies the permanent by \(\exp(-\sum x_i-\sum y_j)\). The next lemma bounds this total cost and also gives the Frobenius displacement needed to compare Gaussian factors.

\paperstatement{Lemma 6.2 (Displacement and capacity).}{6.2}{thm-6.2} Under the preceding assumptions, suppose in addition that \(X\ge0\). Write \(\theta_X=\sum_i x_i+\sum_j y_j\) for the total scaling potential. With constants depending only on the fixed density and gap parameters,
\[\|(x,y)\|_2\le K\|(\alpha,\beta)\|_2,\qquad \|B-X\|_F\le\frac K{\sqrt p}\|(\alpha,\beta)\|_2,\]
\nopagebreak[3]\[0\le\theta_X\le K\|(\alpha,\beta)\|_2^2.\]

For a nonnegative matrix \(X\) of arbitrary positive mass, suppose that \(\widetilde X=pX/\mathfrak m(X)\) satisfies \paperref{thm-6.1}{Lemma 6.1}. Apply the preceding conclusions to \(\widetilde X\). If \(B\) is its scaling, the exact identities are
\[\theta_X=\theta_{\widetilde X}+p\log\frac p{\mathfrak m(X)},\qquad \operatorname{per}X=e^{-\theta_X}\operatorname{per}B.\]

In particular, if \(\kappa=\mathfrak m(X)-p\), then \(\operatorname{per}X\le e^\kappa\operatorname{per}B\). The potential \(\theta_X\) itself need not be nonnegative when the mass is not \(p\).

\textbf{Proof.} The convex potential is
\[\Phi_X(x,y)=\sum_{i,j}X_{ij}e^{x_i+y_j}-\sum_i x_i-\sum_j y_j.\]

Its gradient at the origin is \(g=(\alpha,\beta)\), and the balancing point is a minimum. If \(R=\|(x,y)\|_\infty\), its Hessian along \(tz\), \(0\le t\le1\), is bounded below by \(e^{-2R}\) times the Hessian at \(B\). On the gauge space the latter has least eigenvalue at least \(1-\|B-P_p\|_{\rm op}\). Thus the whole segment has a fixed lower bound \(\lambda>0\). Integrating the gradient gives \(\|z\|_2\le\lambda^{-1}\|g\|_2\). The dense bound on \(X\) and \(|e^{x_i+y_j}-1|\le K|x_i+y_j|\) then give the Frobenius estimate.

Since \(\Phi_X(0)=p\) and \(\Phi_X(z)=p-\theta_X\), convexity gives \(\theta_X\ge0\). Strong convexity gives \(\theta_X\le\|g\|_2^2/(2\lambda)\). The nonunit-mass identity follows by absorbing the scalar \(p/\mathfrak m(X)\) into the two diagonal potentials. Finally, \(p\log(\mathfrak m(X)/p)\le\kappa\) proves the permanent upper bound.

\Needspace{7\baselineskip}
\subsection{Gaussian deletion and centering}\label{sec-6.2}

Deleting \(t\) rows from a dense normalized matrix removes a total squared row norm of order \(t/n\). Applying the log-determinant derivative to the corresponding Gram-matrix loss preserves this order. We then delete columns using the other Gram matrix. This gives a linear deletion error even when the retained matrix is nonprincipal.

For a real matrix \(Z\) with \(\|Z\|_{\rm op}<1\), define
\[\mathcal G(Z)=\det(I-Z^{\mathsf T}Z)^{-1/2}.\]

The same value is obtained with \(ZZ^{\mathsf T}\); padding by zero rows or columns leaves it unchanged.

\paperstatement{Lemma 6.3 (Gram deletion and stability).}{6.3}{thm-6.3} Suppose \(\|Z\|_{\rm op}\le q_*<1\). Delete row set \(I\) and column set \(J\), with remaining sets \(R,T\). Then
\[0\le\log\mathcal G(Z)-\log\mathcal G(Z[R,T])\le\frac{\sum_{i\in I}\|Z_{i,\cdot}\|_2^2+\sum_{j\in J}\|Z_{\cdot,j}\|_2^2}{2(1-q_*^2)}.\]

For a square remaining matrix \(W\) of order \(m\ge1\), put \(u=\mathbf 1_m/\sqrt m\). Then
\[0\le\log\mathcal G(W)-\log\mathcal G(\Pi_mW\Pi_m)\le\frac{\|W^{\mathsf T}u\|_2^2+\|Wu\|_2^2}{2(1-q_*^2)}.\]

For any two real matrices of the same dimensions with operator norm at most \(q_*\),
\[|\log\mathcal G(U)-\log\mathcal G(V)|\le\frac{\|U\|_F+\|V\|_F}{2(1-q_*^2)}\|U-V\|_F.\]

\textbf{Proof.} Let \(P_R\) be the diagonal coordinate projection onto the retained rows. Deleting rows decreases the column Gram matrix by the positive semidefinite matrix \(Z^{\mathsf T}(I-P_R)Z\), whose trace is the sum of the deleted row norms squared. For \(F(H)=-\tfrac12\log\det(I-H)\) and a positive semidefinite direction \(D\),
\[DF(H)[D]=\tfrac12\operatorname{tr}((I-H)^{-1}D)\le\frac{\operatorname{tr}D}{2(1-q_*^2)}.\]

Integrate this bound and then delete columns using the row Gram matrix. Projection off each all-one direction gives the centering bound in the same way. For the last assertion, integrate along the segment between the two Gram matrices and use
\[\|U^{\mathsf T}U-V^{\mathsf T}V\|_*\le(\|U\|_F+\|V\|_F)\|U-V\|_F,\]

where \(\|\cdot\|_*\) is the trace norm. The Gram segment retains the operator gap.

Now let \(S\) be a tournament sign matrix of order \(n\ge4\), \(s=S\mathbf 1_n\), and \(\tau=\|s\|_2^2/(n-1)^2\). Put
\[Y=\frac{S-I}{n-1},\qquad D_n(S)=\det(I-S^{\mathsf T}S/n^2)^{-1/2}.\]

Skew symmetry and paired singular values imply
\[\|Y\|_{\rm op}^2\le q_Y^2=\frac{n(n-1)/2+1}{(n-1)^2}<1.\]

Delete any \(t\) rows and any \(t\) columns, put \(m=n-t\), and set \(Z=\Pi_mY[R,T]\Pi_m\). Each row and column of \(Y\) has squared Euclidean norm \(n/(n-1)^2\). Moreover,
\[\|Y\mathbf 1_n\|_2^2=\|Y^{\mathsf T}\mathbf 1_n\|_2^2=\tau+\frac n{(n-1)^2},\]
\nopagebreak[3]\[\|Y\mathbf 1_J\|_2^2,\ \|Y^{\mathsf T}\mathbf 1_I\|_2^2\le\frac{t^2n}{(n-1)^2}.\]

The preceding lemma therefore gives
\[0\le\log\mathcal G(Y)-\log\mathcal G(Z)\le\frac{t n/(n-1)^2+2[\tau+n(1+t^2)/(n-1)^2]/m}{1-q_Y^2}.\]

The full Gram identity is \(Y^{\mathsf T}Y=(S^{\mathsf T}S+I)/(n-1)^2\). Its difference from \(S^{\mathsf T}S/n^2\) is positive semidefinite and has trace
\[\mathcal T_n=\frac{2n-1}{n(n-1)}+\frac n{(n-1)^2}.\]

Hence \(0\le\log\mathcal G(Y)-\log D_n(S)\le43/(4n)\). If \(m\ge n/2\), the first loss is at most \((24t+18\tau+8)/n\). Both are nonnegative losses from the same full kernel, so
\[|\log\mathcal G(Z)-\log D_n(S)|\le\frac{24t+18\tau+11}{n}.\]

For completeness, \(1-q_Y^2=n(n-3)/(2(n-1)^2)\). Substitution proves the stated constants directly, using \(t\le n/2\). The estimate applies to independent row and column sets \(I,J\); skew symmetry is used only for the original full matrix \(S\).

\Needspace{7\baselineskip}
\subsection{Preparing the tournament submatrices}\label{sec-6.3}

The row and column errors of \(C=2A/(n-1)\) have opposite signs. Multiplying row \(i\) by \((1+a_i)^{-1}\) corrects its original row sum, while multiplying column \(i\) by \((1-a_i)^{-1}\) corrects its original column sum. Applying both corrections leaves a residual marginal error, estimated below. Their main advantage is that the total mass has no first-order change and their restored product is the score penalty \(\Gamma\).

Fix \(0\le a_0<1\) and finite positive parameters \(A_0,B_0\). In this subsection suppose
\[a_i=\frac{s_i}{n-1},\qquad \max_i|a_i|\le a_0,\qquad \tau=\sum_i a_i^2\le A_0\log n,\qquad t\le B_0\log n.\]

Define
\[C=\frac{J-I+S}{n-1}=\frac{2A}{n-1},\qquad \ell_i=\frac1{1+a_i},\qquad r_i=\frac1{1-a_i},\qquad g_i=\frac1{1-a_i^2}.\]

Set \(C'=\operatorname{diag}(\ell)C\operatorname{diag}(r)\), \(F=C'-C\), and \(\nu=\mathfrak m(C')-n\). To compute its mass, write \(v_i=a_i/(1-a_i^2)\) and \(w_i=a_i^2/(1-a_i^2)\), with \(V=\sum_i v_i\) and \(W=\sum_i w_i\). Expanding \(\ell=\mathbf 1-v+w\) and \(r=\mathbf 1+v+w\) gives the exact identity
\[\nu=\frac{W^2-V^2+W+2w^{\mathsf T}Sv}{n-1}.\]

Indeed, \(\mathbf 1^{\mathsf T}Sv=-(n-1)W\) cancels the first-order mass change. The fixed score bound implies \(W\le K\tau\), \(|V|\le K\tau\), and \(\|v\|_2\le K\sqrt\tau\). Estimating \(Sv\) row by row gives
\[|\nu|\le K\mathcal M_\tau,\qquad \mathcal M_\tau=\frac{\tau+\tau^2}{n}+\frac{\tau^{3/2}}{\sqrt n}=o(1).\]

The marginal identities are
\[C'\mathbf 1_n-\mathbf 1_n=\operatorname{diag}(\ell)C(r-\mathbf 1_n),\]
\nopagebreak[3]\[(C')^{\mathsf T}\mathbf 1_n-\mathbf 1_n=\operatorname{diag}(r)C^{\mathsf T}(\ell-\mathbf 1_n).\]

Since each row and column of \(C\) has Euclidean norm \(O(n^{-1/2})\) and \(\|C\|_{\rm op}=O(1)\), the infinity marginal error is \(O(\sqrt{\tau/n})\), the Euclidean marginal error is \(O(\sqrt\tau)\), and direct entrywise estimation gives
\[\|F\|_F\le K\sqrt{\tau/n}.\]

Normalize \(\widehat C=nC'/\mathfrak m(C')\). Write its row and column errors as \(e_i,f_j\), so \(\sum e_i=\sum f_j=0\). For arbitrary deleted sets \(I,J\) of size \(t\), write \(R=[n]\setminus I\), \(T=[n]\setminus J\) and \(m=n-t\). Let
\[X=\frac n m\widehat C[R,T],\qquad \kappa=\mathfrak m(X)-m,\qquad \widetilde X=\frac m{\mathfrak m(X)}X.\]

The exact deletion mass identity is
\[\kappa=\frac n m\left[-\sum_{i\in I}e_i-\sum_{j\in J}f_j+\sum_{i\in I,j\in J}\widehat C_{ij}-\frac{t^2}{n}\right].\]

For each remaining row, the exact marginal error before the last normalization is
\[(X\mathbf 1_m)_i-1=\frac n m\left[e_i+\frac t n-\sum_{j\in J}\widehat C_{ij}\right],\]

and there is the corresponding column formula. Consequently,
\[|\kappa|\le K\left[t\sqrt{\tau/n}+t^2/n\right],\]
\nopagebreak[3]\[\varepsilon(\widetilde X)\le K(\sqrt{\tau/n}+t/n),\qquad \|g(\widetilde X)\|_2\le K(\sqrt\tau+t/\sqrt n).\]

Here \(g(\widetilde X)\) is the concatenated row and column marginal error, not the vertex weight \(g_i\). All entries of \(\widetilde X\) are nonnegative and at most \(K/m\).

To apply \paperref{thm-6.1}{Lemma 6.1}, we now check the centered kernel of \(\widetilde X\). With
\[\eta=\frac{n^2}{\mathfrak m(C')\mathfrak m(X)}=1+O((t+1)/n),\]

it is
\[\Pi_m\widetilde X\Pi_m=\eta\left[Z+\Pi_mF[R,T]\Pi_m\right].\]

Since \(\mathfrak m(C')=n+o(1)\) and \(\mathfrak m(X)=m+o(1)\), the leading value of \(\eta\) is \(n/m\). Compression preserves the operator norm bound for \(Y\), and the displayed Frobenius estimate controls the second term. Thus the centered operator norm is at most \(1/\sqrt2+o(1)\), uniformly under the fixed parameters. Its entries are \(O(1/m)\), and the marginal infinity error tends to zero. The local scaling lemma applies for all sufficiently large \(n\) and constructs a genuine doubly stochastic matrix \(B_X\).

The Euclidean displacement lemma yields
\[\|B_X-\widetilde X\|_F\le K(\sqrt{\tau/n}+t/n).\]

Combining this with the sharper centered bound, the centered norm of \(B_X\) is still \(1/\sqrt2+o(1)\); for example, it is eventually at most \(4/5\). The entries remain bounded by \(K/m\).

Finally, apply the Gaussian deletion estimate to \(Z\), then the logarithmic Lipschitz estimate to the preconditioning, normalization, and true-scaling perturbations. Since \(B_X-P_m=\Pi_mB_X\Pi_m\), this gives
\[|\log\mathcal G(B_X-P_m)-\log D_n(S)|\le K\left[\sqrt{\tau/n}+(t+1)/n\right].\]

The term \(\tau/n\) from centering is absorbed by \(\sqrt{\tau/n}\). The deletion part of this estimate is linear in \(t/n\), as needed when it is summed against the subset weights.

\Needspace{7\baselineskip}
\subsection{Proof of the nonprincipal permanent bound}\label{sec-6.4}

We prove \paperref{thm-3.1}{Lemma 3.1}, using \(n\) for the full core order \(N\) and \(a_0=b\). The preceding construction provides the scaling matrix \(B_X\). It remains to restore every diagonal and scalar factor and compare the resulting Gaussian determinant with \(D_n(S)\).

\textbf{Proof.} The true scaling just constructed satisfies its density and centered-gap hypotheses, so \paperref{thm-5.1}{Theorem 5.1} gives
\[\operatorname{per}B_X=\frac{m!}{m^m}\mathcal G(B_X-P_m)(1+O(1/m)).\]

The row and column factors restore exactly as
\[\prod_{i\in R}(1+a_i)\prod_{j\in T}(1-a_j)=\Gamma\prod_{i\in I}\ell_i\prod_{j\in J}r_j.\]

Thus deleted rows restore \(\ell\) factors and deleted columns restore \(r\) factors. The complete scalar identity is
\[\operatorname{per}A[R,T]=\Gamma\prod_{i\in I}\ell_i\prod_{j\in J}r_j\left[\frac{m(n-1)\mathfrak m(C')}{2n^2}\right]^m e^{-\theta_X}\operatorname{per}B_X.\]

After restoring \(m!/m^m\), the remaining scalar is
\[\frac{m!}{2^m}(1-1/n)^m\left[\frac{\mathfrak m(C')}{n}\right]^m.\]

Its logarithm differs from that of \(e^{-1}m!/2^m\) by \(O((t+1)/n+|\nu|)\). The capacity bound gives \(e^{-\theta_X}\le e^\kappa\). Substituting the mass, Gaussian, and zero-order errors proves the claimed exponent. All constants are uniform over the tournament and the two deletion sets, proving \paperref{thm-3.1}{Lemma 3.1}. \(\square\)

\Needspace{7\baselineskip}
\subsection{The small-score two-sided estimate}\label{sec-6.5}

When every score is \(o(\sqrt n)\), the original margins are already close enough to one for direct local scaling. The capacity estimate then controls the restoration on both sides, giving the approximation used in \paperref{thm-4.1}{Theorem 4.1}.

\paperstatement{Lemma 6.4 (Small-score permanent approximation).}{6.4}{thm-6.4} Let \(d=\|S\mathbf 1_n\|_\infty=o(\sqrt n)\). Fix \(0<B_0<\infty\), delete any \(t\le B_0\log n\) rows and any \(t\) columns, and put \(m=n-t\). Uniformly over these choices,
\[\operatorname{per}A[R,T]=e^{-1}D_n(S)\frac{m!}{2^m}\left[1+O_{B_0}\left(\frac{(d+t+1)^2}{n}\right)\right].\]

For fixed \(B_0\), the implied constant is uniform once \((d+t+1)^2/n\) is sufficiently small. In the stated regime this quantity tends to zero.

\textbf{Proof.} No paired preconditioning is needed. Set \(C=2A/(n-1)\) and \(X=(n/m)C[R,T]\). The row and column errors of the full matrix \(C\) are \(s_i/(n-1)\) and \(-s_i/(n-1)\), respectively. The exact mass and marginal formulas from \paperref{sec-6.3}{Section 6.3} therefore give
\[|\kappa|=O(t(d+t)/n),\qquad \varepsilon(X)=O((d+t)/n),\qquad \|g(X)\|_2=O((d+t)/\sqrt n).\]

Normalize \(\widetilde X=mX/\mathfrak m(X)\). Its density and centered gap are fixed, and the local scaling lemma applies. The Euclidean displacement and capacity estimates give
\[\|B_X-\widetilde X\|_F=O((d+t)/n),\qquad 0\le\theta_{\widetilde X}=O((d+t)^2/n).\]

The nonunit-mass correction then implies \(|\theta_X|=O((d+t)^2/n)\). Also \(\tau\le nd^2/(n-1)^2\). The finite deletion estimate, the normalization factor, and the true-scaling perturbation show
\[|\log\mathcal G(B_X-P_m)-\log D_n(S)|=O((d+t+1)^2/n).\]

The exact restoration, now without score products, is
\[\operatorname{per}A[R,T]=\left[\frac{m(n-1)}{2n}\right]^m e^{-\theta_X}\operatorname{per}B_X.\]

Apply the uniform permanent theorem and use \((1-1/n)^m=e^{-1}\exp(O((t+1)/n))\). All logarithmic errors are \(O((d+t+1)^2/n)=o(1)\), so exponentiating them proves the stated relative approximation. This proves the uniform short-minor approximation used in \paperref{thm-4.1}{Theorem 4.1}. \(\square\)

\Needspace{9\baselineskip}
\section{Further questions}\label{sec-7}

The first question is whether the upper constant can be reduced to \(L\). The spectral argument uses only the largest normalized squared frequency and their total mass. Additional restrictions on spectra realizable by tournament matrices could improve this relaxation; in particular, proving \(\rho_n(S)\le r_n\) for every tournament, with \(r_n\) as in \paperref{sec-4.2}{Section 4.2}, would give \(P(n)=(L+O(n^{-1}))\mu_n\) through the same reduction.

The proof does not evaluate a usable pair of constants \(K,n_0\). Doing so would require tracking the uniform scaling thresholds and the error budgets through all three tournament classes. The \(O(n^{-1})\) approximation also leaves open the first correction coefficient for carousel path counts and the structure of finite-order maximizers.

\appendix

\Needspace{9\baselineskip}
\section{Formalization and reproducibility}\label{sec-A}

\Needspace{7\baselineskip}
\subsection{Formal statements and correspondence}\label{sec-A.1}

The \href{https://github.com/LStar404/tournament-hamiltonian-paths/tree/8ea3fcffcc12b6a06294ba7559885b439eda3cac/formalization}{companion Lean project} uses the same finite objects as this paper: a tournament is a loopless orientation of every pair of distinct labelled vertices; a Hamiltonian path is a vertex permutation whose consecutive arcs point forward; and \(P(n)\) is the maximum of the resulting count. Lean and Mathlib are pinned to version 4.34.1.

The declaration \texttt{Tournament\allowbreak{}Hamiltonian.\allowbreak{}main\allowbreak{}Bound : Tournament\allowbreak{}Hamiltonian.\allowbreak{}Main\allowbreak{}Bound} proves \paperref{thm-1.1}{Theorem 1.1} with the exact constants \(L\) and \(C_*\). Its conclusion chooses one \(K\ge0\) and one \(n_0\ge2\) before quantifying over all \(n\ge n_0\). The final theorem has no unproved permanent, scaling, activity or path-count premise. The separate declaration \texttt{small\_\allowbreak{}score\_\allowbreak{}path\allowbreak{}Count\_\allowbreak{}spectral\_\allowbreak{}approximation\_\allowbreak{}eventually} proves \paperref{thm-4.1}{Theorem 4.1}: its absolute error constant is independent of the score envelope \(d(n)\), while its eventual threshold may depend on that envelope.

The formal development proves the principal analytic inputs as well as these conclusions. Some intermediate arguments differ from the presentation above. The permanent proof uses a geometric coefficient-moment bound in place of the displayed derivative constant \(K_G\). The high-variance and exceptional-vertex cases use weaker sufficient estimates; for the latter the formal bound is \((10/3)4^{-f}(1+o(1))<1\) when \(f\ge1\). The weighted determinant comparison uses positive principal-minor moments. Small-score restoration reuses paired preconditioning, and local scaling uses equivalent contraction estimates with a different fixed gap. The formal short-subset convention differs by one boundary layer, and the polarization and Gaussian-moment tools have finite-dimensional implementations. These choices prove the same stated rates and final bounds. The accompanying {result-by-result correspondence\footnote{Source archive: \path|materials/verification/manuscript_alignment_20261011.md|}} records the intermediate differences.

\Needspace{7\baselineskip}
\subsection{Verification records and auxiliary computations}\label{sec-A.2}

The repository\textquotesingle s \href{https://github.com/LStar404/tournament-hamiltonian-paths/blob/8ea3fcffcc12b6a06294ba7559885b439eda3cac/formalization/VERIFICATION.md}{verification record} reports a successful complete run on 9 October 2026. The command \texttt{python verify\_\allowbreak{}lean.\allowbreak{}py {-}{-}require{-}main} builds the project, checks the actual \texttt{Main\allowbreak{}Bound} declaration, checks project-source import coverage, and audits transitive theorem axioms against \texttt{propext}, \texttt{Classical.\allowbreak{}choice} and \texttt{Quot.\allowbreak{}sound}. The recorded audit contains 3674 theorem constants, including generated lemmas. All 246 stored Lean-source and configuration hashes match the sources inspected for this revision. These are the recorded build results and a source-correspondence check, rather than a new Lean build performed for the editorial revision.

Exact finite computations check normalizations, coefficient identities, deletion factors and the four-block expansion. They are separate from the all-order proof. The archived rational calculations give
\[2.855957892565113<L<2.855957892565114,\]
\nopagebreak[3]\[2.857401177672316<C_*<2.857401177672317,\]

and place \((C_*-L)/L\) between \(0.000505359379058\) and \(0.000505359379059\). The upper-constant enclosure is also proved in Lean; the other decimal enclosures are recorded in the rational computational certificates.

Xiangyu Ye contributed the companion formalization, submitted in \href{https://github.com/LStar404/tournament-hamiltonian-paths/pull/1}{pull request 1}. AI tools assisted proof reconstruction, cross-checking, translation and exact-arithmetic diagnostics.

\Needspace{6\baselineskip}\section*{References}
\phantomsection\addcontentsline{toc}{section}{References}\label{sec-references}\small

\hypertarget{bib-1}{}\phantomsection\label{bib-1}\noindent {[}1{]} John Irving and Mohamed Omar. Revisiting the Rédei-Berge Symmetric Functions via Matrix Algebra. The Electronic Journal of Combinatorics 32(4) (2025), Paper P4.43, 23 pp. \href{https://doi.org/10.37236/13841}{\nolinkurl{DOI: 10.37236/13841}}.

\hypertarget{bib-2}{}\phantomsection\label{bib-2}\noindent {[}2{]} Noga Alon. The Maximum Number of Hamiltonian Paths in Tournaments. Combinatorica 10(4) (1990), 319--324. \href{https://doi.org/10.1007/BF02128667}{\nolinkurl{DOI: 10.1007/BF02128667}}.

\hypertarget{bib-3}{}\phantomsection\label{bib-3}\noindent {[}3{]} Yanjun Han and Jonathan Niles-Weed. Approximate independence of permutation mixtures. \href{https://arxiv.org/html/2408.09341v2}{arXiv:2408.09341v2} (2024).

\hypertarget{bib-4}{}\phantomsection\label{bib-4}\noindent {[}4{]} Bo Deng, Xueliang Li, Bryan Shader and Wasin So. On the Maximum Skew Spectral Radius and Minimum Skew Energy of Tournaments. Linear and Multilinear Algebra 66(7) (2018), 1434--1441. \href{https://doi.org/10.1080/03081087.2017.1357676}{\nolinkurl{DOI: 10.1080/03081087.2017.1357676}}.

\hypertarget{bib-5}{}\phantomsection\label{bib-5}\noindent {[}5{]} Peter McCullagh. An asymptotic approximation for the permanent of a doubly stochastic matrix. Journal of Statistical Computation and Simulation 84(2) (2014), 404--414. \href{https://doi.org/10.1080/00949655.2012.712122}{\nolinkurl{DOI: 10.1080/00949655.2012.712122}}. \href{https://arxiv.org/abs/1205.5723}{arXiv:1205.5723}.

\hypertarget{bib-6}{}\phantomsection\label{bib-6}\noindent {[}6{]} N. C. Wormald. Tournaments with many Hamilton cycles. \href{https://users.monash.edu.au/~nwormald/papers/hamtourn.pdf}{Undated preprint}, 20 pp. Accessed 10 October 2026.

\hypertarget{bib-7}{}\phantomsection\label{bib-7}\noindent {[}7{]} Ehud Friedgut and Jeff Kahn. On the Number of Hamiltonian Cycles in a Tournament. Combinatorics, Probability and Computing 14(5--6) (2005), 769--781. \href{https://doi.org/10.1017/S0963548305006863}{\nolinkurl{DOI: 10.1017/S0963548305006863}}.

\hypertarget{bib-8}{}\phantomsection\label{bib-8}\noindent {[}8{]} Eric Li. The Godsil--McKay Asymptotic for Latin Rectangles in the Sublinear Range of Erdős Problem 725. \href{https://arxiv.org/html/2608.01671v1}{arXiv:2608.01671v1} (2026).
\end{document}
